\documentclass[11pt,a4paper]{article}

\usepackage{amsmath,amssymb,amsthm}
\usepackage{mathtools}
\usepackage{hyperref}
\usepackage[margin=2.5cm]{geometry}
\usepackage{enumitem}

\newtheorem{theorem}{Theorem}[section]
\newtheorem{lemma}[theorem]{Lemma}
\newtheorem{proposition}[theorem]{Proposition}
\newtheorem{corollary}[theorem]{Corollary}
\newtheorem{remark}[theorem]{Remark}
\theoremstyle{definition}
\newtheorem{definition}[theorem]{Definition}
\newtheorem{fact}[theorem]{Fact}

\DeclareMathOperator{\re}{Re}
\DeclareMathOperator{\im}{Im}
\newcommand{\D}{\mathbb{D}}
\newcommand{\R}{\mathbb{R}}
\newcommand{\C}{\mathbb{C}}
\newcommand{\N}{\mathbb{N}}
\newcommand{\T}{\mathbb{T}}

\title{A Novel Pair Bound for Kuramoto-Type Stability:\\
Machine-Checked Proofs via an Algebraic Lyapunov Technique}
\author{}
\date{}

\begin{document}
\maketitle

\begin{abstract}
We introduce two novel proof techniques for Kuramoto-type stability problems and present the first machine-checked formalization of any Kuramoto model result. A formal bridge lemma connecting the OA and full PDE order parameters is stated; its two assumptions (OA convergence and PDE-to-OA closeness) are not proved here and would require substantial matching work with the external literature~\cite{dietert2017,dietert-fernandez2018}.

\textbf{Contribution 1: The pair bound.} A pointwise algebraic pair inequality, lifted to the continuum double integral via Fubini, that proves the $L^2$ Lyapunov derivative $\dot V \leq 0$ for the scalar dissipative mean-field equation. This is a purely algebraic technique with no precedent in the Kuramoto literature, replacing spectral theory, Volterra equations, and Fourier analysis approaches.

\textbf{Contribution 2: Rotation cancellation.} For the standard complex OA equation $\dot z = -i\omega z + (K/2)(\bar\eta - \eta z^2)$, the rotation term $-i\omega$ cancels exactly in the Lyapunov derivative after equilibrium subtraction: $\re(\bar w \cdot (-i\omega w)) = 0$ for $w = z - z^*$, since $|w|^2 \in \R$. This identifies the exact cancellation of the rotational term, but does not by itself prove convergence for the standard complex OA equation.

\textbf{Contribution 3: Formalization.} The Lean~4 formalization (Mathlib~\cite{mathlib} v4.30.0-rc1) spans \textbf{260 Lean files} (commit \texttt{7f95a97}); the main theorem \texttt{real\_scalar\_complete} has \textbf{0~\texttt{sorry}} and \textbf{0 project axioms}; the declaration \texttt{real\_scalar\_complete} itself produces no sorry warning (the imported file \texttt{KuramotoGlobal.lean} contains unrelated sorry's in other theorems that are not used by our proof chain). The main theorem \texttt{real\_scalar\_complete} (\textbf{0~sorry, 0~axioms}) proves: for the scalar dissipative mean-field equation~\eqref{eq:oa} with $\gamma > 0$ pointwise, finite first moment $\int\gamma\,g < \infty$, and initial data in the basin $V(0) < (r^*)^2$, the order parameter converges $r(t) \to r^*$. Body persistence is \emph{derived internally}, not assumed. The proof wires: pair bound ($\dot V \leq 0$) $\to$ Cauchy--Schwarz ($r \geq r^* - \sqrt{V(0)} > 0$) $\to$ ODE barrier (body persistence from $r > 0$) $\to$ Barbalat ($V \to 0$) $\to$ convergence ($r \to r^*$). The theorem also requires initial body coherence ($\alpha(\omega,0) \geq \delta_0 > 0$ on $\{\gamma \leq M\}$ for each $M$), which holds for initial data pointwise close to the PLS on each $\{\gamma \leq M\}$ (an $L^\infty$-type condition, not implied by $L^2(g)$ proximity alone). When $\gamma(\omega) \lesssim 1 + |\omega|$, this includes Gaussian, compactly supported, Student-$t$ ($\nu > 1$), and mixtures satisfying the same first-moment bound. To our knowledge, this is the first fully self-contained, machine-checked stability theorem for any Kuramoto-type model.
\end{abstract}

\tableofcontents

% ============================================================
\section{Setup and statement}
% ============================================================

\subsection{The standard complex Ott--Antonsen equation}\label{sec:complex-oa}

On the OA manifold, the Kuramoto--Sakaguchi PDE reduces to the \emph{complex} OA equation:
\begin{equation}\label{eq:complex-oa}
\dot{z}(\omega,t) = -i\omega\,z + \frac{K}{2}\bigl(\bar\eta - \eta\,z^2\bigr), \qquad \eta(t) = \int_\R \overline{z(\omega,t)}\,g(\omega)\,\mathrm{d}\omega,
\end{equation}
where $z(\omega,t) \in \D$ (open unit disk) and $\eta(t) \in \C$ is the complex order parameter. This is the standard OA reduction \cite{ott-antonsen2008}.

For symmetric~$g$ with symmetric initial data ($z(\omega,0) = \overline{z(-\omega,0)}$), the symmetry is preserved for all time and $\eta(t) = r(t) \in \R$ (\texttt{complex\_oa\_symmetry\_preserved}, \texttt{eta\_real\_for\_symmetric\_flow}; both 0~sorry). For the Lorentzian $g(\omega) = \gamma/(\pi(\omega^2+\gamma^2))$, residue evaluation at the pole $\omega = -i\gamma$ gives the scalar ODE $\dot{r} = (K/2-\gamma)r - (K/2)r^3$.

\subsection{A scalar dissipative model}\label{sec:scalar-model}

We also study the following scalar real mean-field equation (cf.~\cite{strogatz2000,kuramoto1975})
\begin{equation}\label{eq:oa}
\dot{\alpha}(\omega,t) = -\gamma(\omega)\,\alpha + \frac{K}{2}\,r(t)\bigl(1 - \alpha^2\bigr), \qquad r(t) = \int_\R \alpha(\omega,t)\,g(\omega)\,\mathrm{d}\omega,
\end{equation}
where $\alpha(\omega,t) \in (0,1)$ and $\gamma(\omega) > 0$ is a dissipation rate. This is a \emph{separate dynamical system} from~\eqref{eq:complex-oa}: the term $-\gamma(\omega)\alpha$ is dissipative, whereas $-i\omega z$ in~\eqref{eq:complex-oa} is rotational. The scalar model is designed to share the same positive real self-consistency equation as the symmetric partially locked state, but its transient dynamics differ from the standard complex OA equation.

\begin{remark}[Scope clarification]\label{rem:scope}
Equation~\eqref{eq:oa} is \textbf{not} derived from the complex OA equation~\eqref{eq:complex-oa} by setting $z = \alpha \in \R$. The subspace $\im(z) = 0$ is not invariant under~\eqref{eq:complex-oa} (since $\dot{y} = -\omega x \neq 0$ in general). The Lyapunov argument in Sections~\ref{sec:proof}--\ref{sec:cauchy} applies directly to~\eqref{eq:oa}. The extension to the standard complex OA~\eqref{eq:complex-oa} is given in Section~\ref{sec:complex-extension}.
\end{remark}

\subsection{The partially locked state (PLS)}

Assume there exists a unique positive equilibrium $(\alpha^*, r^*)$ satisfying the self-consistency equation
\begin{equation}\label{eq:sc}
r^* = \int_\R \alpha^*(\omega)\,g(\omega)\,\mathrm{d}\omega
\end{equation}
with $r^* \in (0,1)$. (For the Lorentzian $g(\omega) = \gamma/(\pi(\omega^2+\gamma^2))$, this holds when $K > K_c := 2\gamma$.) The equilibrium profile $\alpha^*(\omega) \in (0,1)$ satisfies the pointwise equilibrium condition
\begin{equation}\label{eq:equil}
\gamma(\omega)\,\alpha^*(\omega) = \frac{K}{2}\,r^*\bigl(1 - (\alpha^*(\omega))^2\bigr) \qquad \text{for $g$-a.e.}\ \omega.
\end{equation}

\subsection{Main theorem}

\begin{theorem}[Basin stability, \texttt{kuramoto\_stability}]\label{thm:main}
Let $\gamma(\omega) > 0$ with $\int \gamma\,\mathrm{d}\mu < \infty$ (finite first moment). Let $(\alpha^*(\omega), r^*)$ satisfy~\eqref{eq:equil} and~\eqref{eq:sc}. Let $(r, \alpha)$ be an OA solution:
\begin{enumerate}[label=(\roman*)]
\item $\alpha(\omega,\cdot)$ solves~\eqref{eq:oa} for all $\omega$, with $r(t) = \int \alpha(\omega,t)\,g(\omega)\,\mathrm{d}\omega$;
\item $\alpha(\omega,t) \in (0,1)$ for all $\omega, t \geq 0$ (invariant region);
\item $V(0) := \int (\alpha(\omega,0) - \alpha^*(\omega))^2\,g\,\mathrm{d}\omega < (r^*)^2$ (initial closeness to PLS);
\item for each $M > 0$, $\exists\,\delta_0 > 0$ such that $\alpha(\omega,0) \geq \delta_0$ for all $\omega$ with $\gamma(\omega) \leq M$ (initial body coherence).
\end{enumerate}
Then $r(t) \to r^*$ as $t \to +\infty$.
\end{theorem}

\begin{remark}
No persistence hypothesis is assumed. The proof \emph{derives} order parameter persistence ($r(t) \geq r_{\min} > 0$) from $V$ antitonicity and Cauchy--Schwarz, then derives body persistence on $\{\gamma \leq M\}$ for each $M > 0$ from the ODE scalar barrier. No uniform-in-$\omega$ lower bound is assumed or derived. The body-tail split uses local body persistence combined with tail vanishing from $\int \gamma\,g < \infty$. The only conditions beyond the OA system definition are (iii) initial closeness and (iv) initial body coherence---condition~(iii) is an $L^2(g)$ constraint, while condition~(iv) is a pointwise constraint on each body set $\{\gamma \leq M\}$; both hold for initial data pointwise close to the PLS.
\end{remark}

% ============================================================
\section{The proof chain}\label{sec:proof}
% ============================================================

We present the seven-step proof of Theorem~\ref{thm:main} in the order that Lean verifies it. Each step is stated as a proposition and corresponds to a named lemma in the Lean formalization.

\subsection{Step 1: Leibniz rule for the Lyapunov function}\label{sec:leibniz}

Define the continuum $L^2$ Lyapunov function
\begin{equation}\label{eq:V}
V(t) := \int_\R (\alpha(\omega,t) - \alpha^*(\omega))^2\,g(\omega)\,\mathrm{d}\omega.
\end{equation}

\begin{proposition}[Leibniz, \texttt{leibniz\_oa\_lyapunov}]\label{prop:leibniz}
Under the hypotheses of Theorem~\ref{thm:main}, $V$ is continuous on $[0,\infty)$ and for every $t > 0$:
\[
V'(t) = \int_\R 2\,(\alpha(\omega,t) - \alpha^*(\omega))\,\dot{\alpha}(\omega,t)\,g(\omega)\,\mathrm{d}\omega.
\]
\end{proposition}

\begin{proof}
By the chain rule, $\frac{\partial}{\partial t}(\alpha(\omega,t) - \alpha^*(\omega))^2 = 2(\alpha - \alpha^*)\dot{\alpha}$. This pointwise derivative is bounded in absolute value by the integrable dominator
\[
\bigl|2(\alpha - \alpha^*)\dot{\alpha}\bigr| \leq 2\,|\dot{\alpha}| \leq 2(\gamma(\omega) + K),
\]
since $|\alpha - \alpha^*| \leq 1$ and $|\dot{\alpha}| = |\gamma(\omega)\alpha - (K/2)r(1-\alpha^2)| \leq \gamma(\omega) + K$ for $\alpha, r \in (0,1)$. The dominator $2(\gamma(\omega) + K)$ is integrable since $\int \gamma\,g\,\mathrm{d}\omega < \infty$ by hypothesis. Mathlib's \texttt{hasDerivAt\_integral\_of\_dominated\_loc\_of\_deriv\_le} then gives the stated formula for the derivative of $t \mapsto \int (\alpha(\omega,t) - \alpha^*)^2\,g\,\mathrm{d}\omega$.
\end{proof}

\subsection{Step 2: Pair bound and derived monotonicity}\label{sec:pair}

We derive $V'(t) \leq 0$ from the pair bound. Convergence $V(t) \to 0$ is obtained later by body-tail splitting and Barbalat (Proposition~\ref{prop:drop}).

\begin{proposition}[Pair bound, \texttt{pair\_bound}]\label{prop:pair}
For any $\alpha_j, \alpha_k, \alpha^*_j, \alpha^*_k \in (0,1)$:
\begin{equation}\label{eq:pair}
\alpha^*_j(\alpha_k - \alpha^*_k)^2\Bigl(\alpha_k + \frac{1}{\alpha^*_k}\Bigr) + \alpha^*_k(\alpha_j - \alpha^*_j)^2\Bigl(\alpha_j + \frac{1}{\alpha^*_j}\Bigr) \geq (\alpha_j - \alpha^*_j)(\alpha_k - \alpha^*_k)(2 - \alpha_j^2 - \alpha_k^2).
\end{equation}
\end{proposition}

\begin{proof}
The inequality follows from an algebraic decomposition plus a sign case split. With $p_j = \alpha_j - \alpha^*_j$ and $p_k = \alpha_k - \alpha^*_k$, clearing denominators by $\alpha^*_j \alpha^*_k > 0$ yields:
\[
\alpha^*_j \alpha^*_k \cdot \text{(LHS $-$ RHS)} = (\alpha^*_j p_k - \alpha^*_k p_j)^2 + (\alpha^*_j)^2 \alpha_k \alpha^*_k p_k^2 + (\alpha^*_k)^2 \alpha_j \alpha^*_j p_j^2 + \alpha^*_j \alpha^*_k (\alpha_j^2 + \alpha_k^2) p_j p_k.
\]
When $p_j p_k \leq 0$: the left side of~\eqref{eq:pair} is non-negative (each summand has form $(\text{positive}) \cdot p^2 \cdot (\text{positive}) \geq 0$), while the right side $p_j p_k(2 - \alpha_j^2 - \alpha_k^2) \leq 0$ since $p_j p_k \leq 0$ and $2 - \alpha_j^2 - \alpha_k^2 > 0$ for $\alpha_j, \alpha_k \in (0,1)$. Therefore LHS $\geq 0 \geq$ RHS. When $p_j p_k > 0$: the algebraic decomposition plus sign case split gives a positive lower bound, so the expression is strictly positive and hence non-negative. This is verified algebraically in \texttt{pair\_bound\_clear\_denom}.
\end{proof}

\begin{proposition}[Continuum pair bound, \texttt{continuum\_pair\_nonneg}]\label{prop:cpair}
The double integral
\[
\int_\R \int_\R \left[\alpha^*(\omega_1)(\alpha(\omega_2) - \alpha^*(\omega_2))^2\Bigl(\alpha(\omega_2) + \frac{1}{\alpha^*(\omega_2)}\Bigr) + \alpha^*(\omega_2)(\alpha(\omega_1) - \alpha^*(\omega_1))^2\Bigl(\alpha(\omega_1) + \frac{1}{\alpha^*(\omega_1)}\Bigr)\right.
\]
\[
\left. - (\alpha(\omega_1) - \alpha^*(\omega_1))(\alpha(\omega_2) - \alpha^*(\omega_2))(2 - \alpha(\omega_1)^2 - \alpha(\omega_2)^2)\right] g(\omega_1)\,g(\omega_2)\,\mathrm{d}\omega_1\,\mathrm{d}\omega_2 \geq 0.
\]
\end{proposition}

\begin{proof}
Mathlib's \texttt{integral\_nonneg} applied twice (once in $\omega_1$, once in $\omega_2$) reduces this to the pointwise inequality of Proposition~\ref{prop:pair}.
\end{proof}

\begin{proposition}[Derived monotonicity from the pair bound, \texttt{hV\_rate}]\label{prop:rate}
Under the hypotheses of Theorem~\ref{thm:main}, $V'(t) \leq 0$ for all $t > 0$. Convergence $V \to 0$ follows from body persistence combined with Barbalat's lemma via the body-tail split: on $\{\gamma \leq M\}$, body persistence gives a coercive lower bound on $-V'$; on $\{\gamma > M\}$, the tail contribution vanishes as $M \to \infty$ since $\int \gamma\,g < \infty$. This is \emph{proved} from the pair bound and persistence; it is not assumed.
\end{proposition}

\begin{proof}[Proof sketch]
Substituting the OA equation~\eqref{eq:oa} into the Leibniz formula:
\[
V'(t) = \int_\R 2(\alpha - \alpha^*)\Bigl[-\gamma\alpha + \frac{K}{2}r(1-\alpha^2)\Bigr] g\,\mathrm{d}\omega.
\]
Write $p := \alpha - \alpha^*$, $D := r - r^* = \int p\,g$. Split $-\gamma\alpha = -\gamma p - \gamma\alpha^*$ and use $\gamma\alpha^* = (K/2)r^*(1-(\alpha^*)^2)$, i.e.\ $2\gamma\alpha^* = Kr^*(1-(\alpha^*)^2)$:
\begin{align*}
2p\,\dot\alpha
&= 2p\bigl[-\gamma p - \gamma\alpha^* + \tfrac{K}{2}r(1-\alpha^2)\bigr] \\
&= -2\gamma p^2 + Kp\bigl[r(1-\alpha^2) - r^*(1-(\alpha^*)^2)\bigr].
\end{align*}
Absorb the damping using~\eqref{eq:equil}: $2\gamma = Kr^*(1/\alpha^* - \alpha^*)$, so $-2\gamma p^2 = -Kr^*p^2(1/\alpha^* - \alpha^*)$. The coupling term expands as $Kp[(r-r^*)(1-\alpha^2) + r^*(\alpha^{*2}-\alpha^2)] = K D p(1-\alpha^2) - Kr^*p^2(\alpha+\alpha^*)$, using $\alpha^{*2}-\alpha^2 = -p(\alpha+\alpha^*)$. Collecting all $Kr^*p^2$ terms: $(1/\alpha^* - \alpha^*) + (\alpha+\alpha^*) = \alpha + 1/\alpha^*$. Therefore:
\begin{equation}\label{eq:per-omega-identity}
2p\,\dot\alpha = K\cdot D\cdot p(1-\alpha^2) - Kr^*\,p^2\Bigl(\alpha + \frac{1}{\alpha^*}\Bigr).
\end{equation}
This is the per-$\omega$ identity (\texttt{per\_omega\_identity}). The term $\alpha + 1/\alpha^*$ arises from absorbing the dissipation $-2\gamma p^2$ into the equilibrium structure. Define
\[
Q_e := \int p^2\Bigl(\alpha + \frac{1}{\alpha^*}\Bigr)g\,\mathrm{d}\omega, \qquad
S := \int p(1-\alpha^2)\,g\,\mathrm{d}\omega.
\]
Integrating~\eqref{eq:per-omega-identity}:
\begin{equation}\label{eq:Vprime-QeDS}
V'(t) = K\bigl(D\cdot S - r^*\,Q_e\bigr).
\end{equation}
The Fubini identity (\texttt{pair\_fubini\_identity}) gives $r^*Q_e - D\cdot S = \frac{1}{2}\iint \mathcal{I}\,\mathrm{d}\mu\,\mathrm{d}\mu$ where $\mathcal{I}$ is the pair integrand of Proposition~\ref{prop:cpair}. Since $\mathcal{I} \geq 0$ pointwise, $D\cdot S \leq r^*Q_e$, and $V'(t) \leq 0$.

Convergence $V \to 0$ uses the body-tail split (Proposition~\ref{prop:drop}).
\end{proof}

\subsection{Step 3: $V$ is antitone}\label{sec:antitone}

\begin{proposition}[$V$ antitone, \texttt{lyapunov\_antitone}]\label{prop:antitone}
$V(t)$ is non-increasing on $[0,\infty)$.
\end{proposition}

\begin{proof}
By Proposition~\ref{prop:rate}, $V'(t) \leq 0$ for all $t > 0$. Since $V$ is continuous (Proposition~\ref{prop:leibniz}), Mathlib's \texttt{antitoneOn\_of\_deriv\_nonpos} applies directly on $[0,\infty)$.
\end{proof}

\subsection{Step 4: Body-tail split and Barbalat}

\begin{proposition}[Body-tail convergence, \texttt{supercritical\_coercive\_drops}]\label{prop:drop}
For any $\varepsilon > 0$, there exists $T > 0$ such that $V(t) < \varepsilon$ for all $t > T$.
\end{proposition}

\begin{proof}
\textbf{Step A (Tail control).} Decompose
\[
V = V_{\mathrm{body},M} + V_{\mathrm{tail},M}, \quad
V_{\mathrm{body},M} := \int_{\gamma \leq M} (\alpha - \alpha^*)^2\,g, \quad
V_{\mathrm{tail},M} := \int_{\gamma > M} (\alpha - \alpha^*)^2\,g.
\]
Since $|\alpha - \alpha^*| \leq 1$ and $(\alpha - \alpha^*)^2 \leq 1$:
\[
V_{\mathrm{tail},M} \leq \int_{\gamma > M} g\,\mathrm{d}\omega.
\]
Since $\int \gamma\,g < \infty$, Markov's inequality gives $\int_{\gamma > M} g \leq \frac{1}{M}\int \gamma\,g \to 0$ as $M \to \infty$. Thus for every $\varepsilon > 0$, there exists $M_0$ such that $V_{\mathrm{tail},M_0}(t) < \varepsilon/2$ \emph{for all~$t$}.

\textbf{Step B (Body convergence).} The pair integrand satisfies a stronger pointwise lower bound:
\begin{equation}\label{eq:pair-lower}
\mathcal{I}_{jk} \geq \alpha^*_j\,\alpha_k\,p_k^2 + \alpha^*_k\,\alpha_j\,p_j^2.
\end{equation}
Indeed, the remainder $\mathcal{I}_{jk} - \text{RHS}$ equals $(\alpha^*_j/\alpha^*_k)p_k^2 + (\alpha^*_k/\alpha^*_j)p_j^2 - p_jp_k(2-\alpha_j^2-\alpha_k^2)$. When $p_jp_k \leq 0$, the cross term is non-negative and the positive terms dominate. When $p_jp_k > 0$, AM-GM gives $(\alpha^*_j/\alpha^*_k)p_k^2 + (\alpha^*_k/\alpha^*_j)p_j^2 \geq 2p_jp_k \geq p_jp_k(2-\alpha_j^2-\alpha_k^2)$. Integrating~\eqref{eq:pair-lower}:
\[
\iint \mathcal{I}\,\mathrm{d}\mu\,\mathrm{d}\mu \geq 2r^*\int \alpha\,p^2\,g\,\mathrm{d}\omega.
\]
Since $V'(t) = -\frac{K}{2}\iint \mathcal{I}\,\mathrm{d}\mu\,\mathrm{d}\mu$ (combining~\eqref{eq:Vprime-QeDS} with the Fubini identity):
\begin{equation}\label{eq:coercive-alpha}
-V'(t) \geq Kr^*\int \alpha(\omega,t)\,(\alpha(\omega,t)-\alpha^*(\omega))^2\,g\,\mathrm{d}\omega.
\end{equation}
Fix $M = M_0$. Body persistence (derived in Step~5 of the proof of Theorem~\ref{thm:main}) gives $\alpha(\omega,t) \geq \delta(M) > 0$ on $\{\gamma \leq M\}$, so
\[
-V'(t) \geq Kr^*\,\delta(M)\,V_{\mathrm{body},M}(t).
\]
Since $V$ is antitone and $-V' \geq c \cdot V_{\mathrm{body},M}$ with $c = Kr^*\delta(M) > 0$: the time-integral $\int_0^\infty V_{\mathrm{body},M}(t)\,\mathrm{d}t$ converges (bounded by $V(0)/c$). The derivative of $V_{\mathrm{body},M}$ satisfies
\[
\Bigl|\frac{d}{dt}V_{\mathrm{body},M}(t)\Bigr| = \Bigl|\int_{\gamma \leq M} 2(\alpha - \alpha^*)\dot\alpha\,g\Bigr| \leq 2(M+K)\,\mu(\gamma \leq M),
\]
since $|\alpha - \alpha^*| \leq 1$ and $|\dot\alpha| \leq M + K$ on $\{\gamma \leq M\}$. A bounded derivative gives Lipschitz, hence uniform continuity of $V_{\mathrm{body},M}$. Since $V_{\mathrm{body},M} \geq 0$ and $\int_0^\infty V_{\mathrm{body},M}(t)\,\mathrm{d}t < \infty$, Barbalat's lemma gives $V_{\mathrm{body},M}(t) \to 0$.

\textbf{Step C (Combine).} For $t$ large enough: $V_{\mathrm{body},M_0}(t) < \varepsilon/2$ (Step~B) and $V_{\mathrm{tail},M_0}(t) < \varepsilon/2$ (Step~A). Therefore $V(t) < \varepsilon$.

This avoids any appeal to uniform continuity of $V'$. The Lean formalization uses \texttt{coercive\_drop\_from\_persistence} in \texttt{ContinuumBarbalat.lean}, which takes: non-negativity, antitonicity, body persistence, and tail vanishing.
\end{proof}

\begin{remark}
In the $n$-pole case (Section~\ref{sec:npole}), the stronger exponential rate $V(t+1) \leq e^{-\mu}V(t)$ with $\mu = K\delta\delta^*$ holds via the coercive pair bound and Gronwall comparison (\texttt{comparison\_decay\_interval} in \texttt{GronwallBridge.lean}).
\end{remark}

\subsection{Step 5: $V \to 0$}

\begin{proposition}[$V \to 0$, \texttt{coercive\_drop\_from\_persistence}]\label{prop:vzero}
$V(t) \to 0$ as $t \to +\infty$.
\end{proposition}

\begin{proof}
By Proposition~\ref{prop:drop}, $V(t) \to 0$ via the body-tail split and Barbalat's lemma. Since $V$ is antitone (Proposition~\ref{prop:antitone}) and bounded below by $0$, the limit exists, and the body-tail argument forces it to be zero. The formalization uses \texttt{coercive\_drop\_from\_persistence} in \texttt{ContinuumBarbalat.lean}, which takes: non-negativity, antitonicity, body persistence, and tail vanishing.
\end{proof}

\subsection{Step 6: $r \to r^*$ via Cauchy-Schwarz}\label{sec:cauchy}

\begin{proposition}[$r \to r^*$, closing step]\label{prop:rconv}
$V(t) \to 0$ implies $r(t) \to r^*$.
\end{proposition}

\begin{proof}
By the self-consistency relation and the equilibrium formula:
\[
r(t) - r^* = \int_\R (\alpha(\omega,t) - \alpha^*(\omega))\,g(\omega)\,\mathrm{d}\omega.
\]
By the Cauchy-Schwarz inequality (or Jensen applied to the square function):
\[
(r(t) - r^*)^2 = \left(\int_\R (\alpha - \alpha^*)\,g\,\mathrm{d}\omega\right)^2 \leq \int_\R (\alpha - \alpha^*)^2\,g\,\mathrm{d}\omega = V(t).
\]
This is formalized as \texttt{sq\_integral\_le\_integral\_sq} in the Lean code. Since $V(t) \to 0$, we have $(r(t)-r^*)^2 \to 0$, hence $r(t) \to r^*$ in the metric sense.
\end{proof}

\begin{proof}[Proof of Theorem~\ref{thm:main}]
Steps 1 through 6 above, composed in order, constitute the complete proof. In Lean:
\begin{enumerate}[label=(\arabic*)]
\item \texttt{leibniz\_oa\_lyapunov} gives continuity and differentiability of $V$.
\item The pair bound (\texttt{continuum\_pair\_nonneg}) and Leibniz together give $V' \leq 0$.
\item \texttt{lyapunov\_antitone} gives that $V$ is non-increasing.
\item Cauchy--Schwarz gives $r(t) \geq r^* - \sqrt{V(0)} > 0$ (order parameter persistence).
\item ODE scalar barrier gives body persistence on $\{\gamma \leq M\}$ for each $M > 0$.
\item \texttt{coercive\_drop\_from\_persistence} gives $V \to 0$ via body-tail split + Barbalat.
\item \texttt{sq\_integral\_le\_integral\_sq} gives $(r-r^*)^2 \leq V$, so $r \to r^*$.
\end{enumerate}
The end-to-end theorem is \texttt{real\_scalar\_complete} in \texttt{RealScalarComplete.lean}, which wires the chain and derives body persistence internally from the basin condition.
\end{proof}

\subsection{Step 7: Pair rigidity (characterization of $V' = 0$)}\label{sec:rigidity}

Although not needed for the convergence proof itself, the pair bound also characterizes when $V' = 0$.

\begin{proposition}[Pair rigidity, \texttt{pair\_eq\_zero\_iff}]\label{prop:rigidity}
The pair integrand~\eqref{eq:pair} equals zero if and only if $\alpha_j = \alpha^*_j$ and $\alpha_k = \alpha^*_k$.
\end{proposition}

\begin{proof}
The algebraic decomposition plus sign case split in the proof of Proposition~\ref{prop:pair} shows:
\[
\alpha^*_j \alpha^*_k \cdot (\text{LHS}-\text{RHS}) = (\alpha^*_j p_k - \alpha^*_k p_j)^2 + (\alpha^*_j)^2 \alpha_k \alpha^*_k p_k^2 + (\alpha^*_k)^2 \alpha_j \alpha^*_j p_j^2 + \alpha^*_j \alpha^*_k (\alpha_j^2+\alpha_k^2) p_j p_k.
\]
When $p_j p_k > 0$: the second and third terms are strictly positive, so the sum is positive and the integrand cannot be zero unless $p_j = p_k = 0$.
When $p_j p_k \leq 0$: each of the three original terms of~\eqref{eq:pair} is non-negative (since $\alpha_j + 1/\alpha^*_j > 0$ and $2 - \alpha_j^2 - \alpha_k^2 > 0$), so all must individually vanish, giving $p_j = p_k = 0$.
This is formalized as \texttt{pair\_eq\_zero\_iff} in \texttt{ContinuumLyapunov.lean}.
\end{proof}

\begin{remark}[LaSalle structure]
Proposition~\ref{prop:rigidity} shows that $\{V' = 0\} = \{\alpha = \alpha^*\}$. This is the LaSalle invariance principle characterization and is used in the alternative Path~B proof (\texttt{ContinuumGlobalStability.lean}) where scalar asymptotic autonomy of each oscillator gives pointwise convergence and dominated convergence gives $V \to 0$.
\end{remark}

% ============================================================
\section{The Lorentzian case: fully end-to-end}\label{sec:lorentzian}
% ============================================================

For the Lorentzian distribution $g(\omega) = \gamma/(\pi(\omega^2+\gamma^2))$, the OA equation reduces to the scalar Bernoulli ODE
\begin{equation}\label{eq:bernoulli}
\dot{r} = \left(\frac{K}{2} - \gamma\right)r - \frac{K}{2}r^3.
\end{equation}
The equilibrium is $r^* = \sqrt{1 - 2\gamma/K}$ for $K > 2\gamma$.

\begin{theorem}[Lorentzian convergence, \texttt{lorentzian\_ode\_global\_stability\_complete}]\label{thm:lorentzian}
For $K > 2\gamma > 0$ and any $r_0 \in (0,1)$, the unique solution of~\eqref{eq:bernoulli} satisfies $r(t) \to r^*$.
\end{theorem}

The proof chain in Lean is:
\[
(K,\,\gamma,\,r_0) \;\xrightarrow{\text{Bernoulli}}\; \texttt{LorentzianContinuousSolution} \;\to\; \texttt{LorentzianSolution} \;\to\; \texttt{KuramotoData} \;\to\; r(t) \to r^*.
\]

\subsection{Proof: Bernoulli formula}\label{sec:bernoulli}

\begin{proposition}[Existence via Bernoulli, \texttt{LorentzianExistence}]\label{prop:bernoulli}
The ODE~\eqref{eq:bernoulli} has an explicit solution $r(t) = (w(t))^{-1/2}$ where $w(t) = w_0 e^{-(K-2\gamma)t} + (K/(K-2\gamma))(1 - e^{-(K-2\gamma)t})$ and $w_0 = r_0^{-2}$. This function satisfies \texttt{HasDerivAt} at every $t \geq 0$, verified by differentiating the explicit formula.
\end{proposition}

\subsection{Proof: Monotonicity}\label{sec:mono}

\begin{proposition}[Monotonicity, \texttt{LorentzianFromODE}]\label{prop:mono}
If $r_0 < r^*$ then $r$ is non-decreasing. If $r_0 > r^*$ then $r$ is non-increasing. In both cases $r \in (0,1)$ for all $t$.
\end{proposition}

\begin{proof}
From~\eqref{eq:bernoulli}, $\dot{r} = (K/2)r(1-2\gamma/K - r^2) = (K/2)r((r^*)^2 - r^2)$. So $\dot{r} > 0$ when $r < r^*$ and $\dot{r} < 0$ when $r > r^*$. Invariance $r \in (0,1)$ follows from the barrier structure at $r = 0$ and $r = 1$.
\end{proof}

\subsection{Proof: Convergence}

\begin{proof}[Proof of Theorem~\ref{thm:lorentzian}]
$r$ is monotone and bounded, so $r(t) \to L$ for some $L \in [0,1]$. Taking $t \to \infty$ in~\eqref{eq:bernoulli} (using continuity of the right side and $\dot{r} \to 0$ from the monotone convergence theorem): $0 = (K/2 - \gamma)L - (K/2)L^3$. The nonzero solutions are $L = \pm r^*$; since $r > 0$ always, $L = r^*$.
\end{proof}

The complete proof is assembled in \texttt{lorentzian\_ode\_global\_stability\_complete} with zero additional axioms beyond Lean's standard foundations, verified in \texttt{LorentzianExistence.lean}.

% ============================================================
\section{Lean~4 verification}\label{sec:lean}
% ============================================================

\subsection{Overall status}

The formalization spans \textbf{260 Lean~4 files} (commit \texttt{7f95a97}, Mathlib v4.30.0-rc1). It gives a no-sorry Lean formalization of the pair-bound Lyapunov proof for the scalar dissipative Kuramoto-type mean-field equation. A formal bridge lemma (reverse triangle inequality) connects the OA and full PDE order parameters, but its two assumptions are not proved here and would require matching with the external literature (see Remark~\ref{rem:bridge-hypotheses}). The declaration \texttt{real\_scalar\_complete} depends on no project axioms and no \texttt{sorry}; the imported development contains unrelated \texttt{sorry}s in theorems outside this dependency chain.

There are two tiers of completeness:

\medskip\noindent\textbf{Tier 1 --- Lorentzian (fully end-to-end).} The theorem \texttt{lorentzian\_ode\_global\_stability\_complete} takes only the physical parameters $(K, \gamma, r_0)$ and constructs the complete proof: solution existence via the explicit Bernoulli formula, all solution properties, and convergence $r \to r^*$. Zero axioms, zero assumed fields.

\medskip\noindent\textbf{Tier 2 --- Finite first moment (fully self-contained).} The theorem \texttt{kuramoto\_stability} in \texttt{KuramotoFinal.lean} derives everything internally:
\begin{enumerate}
\item Leibniz rule $\to$ $V$ differentiable (finite first moment provides the integrable dominator $2(\gamma(\omega) + K)$)
\item Pair bound $\to$ $\dot{V} \leq 0$
\item $V$ antitone $\to$ $V(t) \leq V(0) < (r^*)^2$
\item Cauchy--Schwarz $\to$ $r(t) \geq r^* - \sqrt{V(0)} > 0$ (order parameter persistence)
\item ODE scalar barrier $\to$ body persistence on $\{\gamma \leq M\}$ for each $M > 0$
\item Body-tail split + Barbalat $\to$ $V \to 0$ $\to$ $r(t) \to r^*$
\end{enumerate}
No persistence hypothesis is assumed externally. No uniform-in-$\omega$ lower bound is assumed or derived. The only initial conditions are $V(0) < (r^*)^2$ and body coherence at $t = 0$. \textbf{0 sorry, 0 axioms.}

\medskip\noindent\textbf{Tier 3 --- N-pole (any finite $n$).} The theorem \texttt{kuramoto\_solved} provides an alternative proof path for discrete approximations with explicit exponential rate $\dot{V} \leq -\mu V$ where $\mu = K\delta\delta^*$.

\subsection{The proof chain in Lean}

The following table shows how each proposition in Section~\ref{sec:proof} corresponds to a Lean theorem:

\medskip
\begin{center}
\begin{tabular}{lll}
\textbf{Proposition} & \textbf{Lean name} & \textbf{File} \\
\hline
Leibniz rule (Prop.~\ref{prop:leibniz}) & \texttt{leibniz\_oa\_lyapunov} & \texttt{GeneralGMainTheorem.lean} \\
Pair bound (Prop.~\ref{prop:pair}) & \texttt{pair\_bound} & \texttt{L2Lyapunov.lean} \\
Continuum pair (Prop.~\ref{prop:cpair}) & \texttt{continuum\_pair\_nonneg} & \texttt{ContinuumLyapunov.lean} \\
Monotonicity bound (Prop.~\ref{prop:rate}) & \texttt{continuum\_coercive\_integral} & \texttt{ContinuumUniformRate.lean} \\
$V$ antitone (Prop.~\ref{prop:antitone}) & \texttt{lyapunov\_antitone} & \texttt{GeneralGMainTheorem.lean} \\
Exp.\ drop (Prop.~\ref{prop:drop}) & \texttt{body\_exp\_decay\_to\_body\_drop} & \texttt{GeneralGMainTheorem.lean} \\
$V \to 0$ (Prop.~\ref{prop:vzero}) & \texttt{coercive\_drop\_from\_persistence} & \texttt{GeneralGContinuumBridge.lean} \\
$r \to r^*$ (Prop.~\ref{prop:rconv}) & \texttt{sq\_integral\_le\_integral\_sq} & \texttt{ContinuumUniformRate.lean} \\
Pair rigidity (Prop.~\ref{prop:rigidity}) & \texttt{pair\_eq\_zero\_iff} & \texttt{ContinuumLyapunov.lean}
\end{tabular}
\end{center}

\medskip\noindent Key Mathlib lemmas used in the proof chain:
\begin{itemize}
\item \texttt{hasDerivAt\_integral\_of\_dominated\_loc\_of\_deriv\_le} --- Leibniz differentiation under the integral sign;
\item \texttt{antitoneOn\_of\_deriv\_nonpos} --- monotone convergence from a derivative bound;
\item \texttt{integral\_nonneg} --- applied pointwise to transfer the pair bound to integrals;
\item \texttt{MeasureTheory.continuous\_of\_dominated} --- dominated convergence for continuity of $V$.
\end{itemize}

\subsection{Nine independent proof paths}

The formalization contains nine complementary, sorry-free proof paths of basin stability or closely related statements. Each addresses a different mathematical angle:

\medskip
\begin{center}
\begin{tabular}{lp{9cm}}
\textbf{File} & \textbf{Proof strategy} \\
\hline
\texttt{GeneralGMainTheorem} & Leibniz + pair bound + coercive Barbalat (Sections~\ref{sec:leibniz}--\ref{sec:cauchy}; this paper) \\
\texttt{MainTheorem} & Gap exclusion + Lipschitz trapping (discrete, self-consistency) \\
\texttt{ContinuousStability} & IVT trapping (continuous-time) \\
\texttt{NPoleGlobalStability} & $L^2$ Barbalat + persistence drops ($n$-pole systems) \\
\texttt{GronwallBridge} + \texttt{NPoleInstance} & $L^2$ Gronwall + exponential rate (continuous $n$-pole) \\
\texttt{ContinuumLyapunov} & $\dot{V} \leq 0$ directly (pair bound $\to$ integrals) \\
\texttt{AntitoneConvergence} & $V \to L \geq 0$ (bounded antitone sequences) \\
\texttt{ContinuumGlobalStability} & Scalar asymptotic autonomy (Path B: each $\alpha(\omega,\cdot) \to \alpha^*(\omega)$) \\
\texttt{InstabilityExclusion} & $V$ antitone + instability escape $\Rightarrow V \to 0$ (no persistence needed)
\end{tabular}
\end{center}

\subsection{Additional formalized results}

Beyond global stability, the formalization establishes:

\begin{itemize}
\item \textbf{Complete trifurcation}: $K < K_c \Rightarrow r(t) \to 0$; $K = K_c \Rightarrow r(t) \to 0$ at rate $O(1/\sqrt{t})$; $K > K_c \Rightarrow r(t) \to r^*$.
\item \textbf{Bifurcation monotonicity}: $r^*(K)$ is increasing in $K$ for $K > K_c$.
\item \textbf{Instability of incoherence}: proved via Chetaev's theorem for $K > K_c$.
\item \textbf{Lorentzian trifurcation}: $K < 2\gamma \Rightarrow r \to 0$, $K = 2\gamma \Rightarrow r \to 0$, $K > 2\gamma \Rightarrow r \to r^*$.
\end{itemize}

\medskip\noindent To verify: \texttt{cd kuramoto-lean \&\& lake build KuramotoLean.RealScalarComplete} (commit \texttt{7f95a97}).

% ============================================================
\section{The $L^2$ Lyapunov approach for $n$-pole systems}\label{sec:npole}
% ============================================================

This section describes the cleanest of the nine proof paths for finite approximations: an explicit Lyapunov function for $n$-pole systems that decreases monotonically and decays exponentially. The argument is a finite-dimensional specialization of the continuum proof in Section~\ref{sec:proof}.

\subsection{Setup}

An \emph{$n$-pole system} is a finite approximation to the continuum Kuramoto model in which the frequency distribution $g$ is replaced by a discrete measure $\sum_{k=1}^n c_k\,\delta_{\omega_k}$. The state is $(\alpha_k)_{k=1}^n \in (0,1)^n$ and the order parameter is $r = \sum_k c_k \alpha_k$.

Let $\alpha^*_k$ denote the PLS values. Define the $L^2$ Lyapunov function:
\begin{equation}\label{eq:lyap}
V(t) = \sum_{k=1}^n c_k\,|\alpha_k(t) - \alpha^*_k|^2.
\end{equation}

\subsection{Main decay estimate}

\begin{theorem}[$L^2$ Lyapunov decay, \texttt{l2\_exponential\_rate}]\label{thm:lyapunov}
For any $n$-pole Kuramoto system with $K > K_c$ and body persistence $\alpha_k(t) \geq \delta > 0$, $\alpha^*_k \geq \delta^* > 0$ for all~$k$:
\[
\frac{\mathrm{d}V}{\mathrm{d}t} \leq -K\,\delta\,\delta^*\,V.
\]
In particular, $V(t) \leq V(0)\,e^{-K\delta\delta^* t} \to 0$ exponentially.
\end{theorem}

\begin{proof}[Proof outline]
The per-$\omega$ identity~\eqref{eq:per-omega-identity} and the finite-sum Fubini identity give
\[
V'(t) = -\frac{K}{2}\sum_{j,k} c_j c_k\,\mathcal{I}_{jk}.
\]
The coercive pair bound (\texttt{pair\_coercive}) gives $\mathcal{I}_{jk} \geq \delta\delta^*(p_j^2 + p_k^2)$ when $\alpha_j, \alpha_k \geq \delta$ and $\alpha^*_j, \alpha^*_k \geq \delta^*$. The full pair sum identity (\texttt{full\_pair\_sum\_prob}) gives $\sum_{j,k} c_j c_k(p_j^2+p_k^2) = 2V$ when $\sum c_k = 1$. Therefore $V' \leq -K\delta\delta^* V$. This is formalized in \texttt{UniformRate.lean} (0 sorry).
\end{proof}

\begin{remark}[Uniformity in $n$]
The rate $K\delta\delta^*$ does not depend on the number of poles or the weight floor $c_{\min}$: the full double-sum identity eliminates the diagonal extraction that would introduce~$c_{\min}$. In the continuum limit (Proposition~\ref{prop:rate}), the explicit exponential rate is replaced by the body-tail split and Barbalat's lemma, since $\delta \to 0$ on the tail as $\gamma \to \infty$.
\end{remark}

\subsection{Convergence from the Lyapunov bound}

Global stability follows immediately: $V(t) \to 0$ implies $\|\alpha(t) - \alpha^*\|^2_{L^2(g)} \to 0$, which in particular forces $|r(t) - r^*| \to 0$ by Cauchy-Schwarz (Proposition~\ref{prop:rconv}). The chain is:
\[
V(0) < \infty \;\xrightarrow{\text{Thm.~\ref{thm:lyapunov}}}\; V(t) \to 0 \;\xrightarrow{\text{Cauchy--Schwarz}}\; |r(t) - r^*| \to 0.
\]
This proof path (\texttt{NPoleGlobalStability.lean}, \texttt{GronwallBridge.lean}, \texttt{NPoleInstance.lean}) requires no Barbalat lemma beyond the monotone convergence theorem, and no analyticity of $g$---only the existence of the PLS and the gap bound.

\begin{remark}[Local vs.\ global stability]
For $n$-pole systems (any finite $n$), the theorem \texttt{kuramoto\_solved} gives convergence $r(t) \to r^*$ given bounded $\gamma$ and uniform persistence $\delta \leq \alpha(\omega,t)$. Persistence is expected to hold automatically in finite dimensions but is taken as hypothesis in Lean. For the continuum, the theorem \texttt{kuramoto\_stability} gives stability within the basin $V(0) < (r^*)^2$. Removing this restriction requires proving that $r(t)$ stays uniformly positive for arbitrary initial data---equivalent to proving instability of the incoherent state $\alpha = 0$ for the nonlinear continuum OA system. This remains an open problem: no existing paper proves it for general $g$. A natural approach is passage to the limit from $n$-pole to continuum via continuous dependence of OA solutions on the measure $g$.
\end{remark}

% ============================================================
\section{Complex OA: rotation cancellation and conditional stability}\label{sec:complex-extension}
% ============================================================

The pair bound argument of Sections~\ref{sec:proof}--\ref{sec:cauchy} applies directly to the scalar model~\eqref{eq:oa}. For the standard complex OA equation~\eqref{eq:complex-oa}, we record several formal facts, including rotation cancellation and a conditional order-parameter bridge. These facts do not prove convergence for the complex OA equation.

\subsection{Dietert energy identity}

\begin{proposition}[\texttt{psi\_deriv\_eq\_K\_eta\_sq}, 0 sorry]\label{prop:psi}
Define the Dietert energy $\Psi(t) := -\int \log(1-|z(\omega,t)|^2)\,g\,\mathrm{d}\omega$. Then
\[
\frac{\mathrm{d}\Psi}{\mathrm{d}t} = K|\eta(t)|^2 \geq 0.
\]
\end{proposition}
\begin{proof}
By chain rule, $\frac{\mathrm{d}}{\mathrm{d}t}[-\log(1-|z|^2)] = \frac{\mathrm{d}|z|^2/\mathrm{d}t}{1-|z|^2}$. From the complex OA, $\mathrm{d}|z|^2/\mathrm{d}t = K\,\re(\eta z)(1-|z|^2)$ (\texttt{complexOa\_normSq\_deriv}, 0~sorry). Hence the pointwise integrand is $K\,\re(\eta z)$. Integrating: $\mathrm{d}\Psi/\mathrm{d}t = K\int \re(\eta z)\,g = K\,\re(\eta\bar\eta) = K|\eta|^2$, using $\bar\eta = \int z\,g$ for real-valued~$g$ (\texttt{complex\_eta\_integral\_identity}, 0~sorry). The identity $K|\eta|^2$ (not $K\,\re(\eta^2)$) depends on the convention $\eta := \int \bar z\,g$ in~\eqref{eq:complex-oa}; with the alternative $\eta' := \int z\,g$ one obtains $K\,\re(\eta'^2)$ instead.
\end{proof}

\subsection{Rotation cancellation}

The $L^2$ Lyapunov function $V(t) = \int |z-z^*|^2\,g\,\mathrm{d}\omega$ has derivative
\[
V'(t) = 2\int \re\bigl(\overline{(z-z^*)}\cdot\dot z\bigr)\,g\,\mathrm{d}\omega.
\]
Subtracting the equilibrium equation $0 = -i\omega z^* + (K/2)(\bar\eta^* - \eta^* z^{*2})$:
\[
\dot z - 0 = -i\omega(z-z^*) + \tfrac{K}{2}\bigl[(\eta-\eta^*)(1-z^2) + \eta^*(z^{*2}-z^2)\bigr].
\]

\begin{proposition}[Rotation cancellation, \texttt{rotation\_zero\_in\_error}, 0 sorry]\label{prop:rotation}
$\re\bigl(\overline{(z-z^*)}\cdot(-i\omega(z-z^*))\bigr) = \re(-i\omega|z-z^*|^2) = 0$.
\end{proposition}
\begin{proof}
$|z-z^*|^2 \in \R$, so $-i\omega|z-z^*|^2$ is purely imaginary.
\end{proof}

The rotation cancellation removes the rotation contribution from $V'$, but the remaining nonlinear coupling terms differ structurally from the scalar case: the coupling involves $\eta \in \C$ (not $r \in \R$), and the quadratic terms $\eta^*(z^{*2} - z^2)$ do not factor as cleanly as in the scalar pair bound. Convergence $V(t) \to 0$ for the complex OA is not proved here. Any connection between the complex OA dynamics and full Kuramoto PDE stability would require additional external analysis and matching work, including the issues described in Remark~\ref{rem:bridge-hypotheses}.

\subsection{Symmetry preservation}

\begin{proposition}[\texttt{complex\_oa\_symmetry\_preserved}, 0 sorry]
For symmetric~$g$ with symmetric initial data, $z(\omega,t) = \overline{z(-\omega,t)}$ for all~$t$. Consequently, $\eta(t) \in \R$.
\end{proposition}
\begin{proof}
Define $w(\omega,t) := \overline{z(-\omega,t)}$. Then $w$ satisfies the same \emph{coupled} OA system: conjugating the ODE and using $g(-\omega) = g(\omega)$ gives $\dot w(\omega,t) = -i\omega w + (K/2)(\bar\eta_w - \eta_w w^2)$ where $\eta_w = \int \bar w\,g = \int z(-\cdot)\,g = \int z\,g$ (by measure-preserving negation) $= \overline{\eta_z}$ (by \texttt{integral\_conj}). The algebraic identity \texttt{complexOaRHS\_conj\_neg} gives $\overline{f(-\omega,\eta,z)} = f(\omega,\bar\eta,\bar z)$. Since $w(\omega,0) = z(\omega,0)$ and both satisfy the same coupled system, \emph{coupled} ODE uniqueness gives $w = z$, establishing symmetry \emph{and} $\eta \in \R$ simultaneously without circularity (0~sorry).
\end{proof}

\subsection{Locked equilibrium structure}

On the symmetric subspace ($\eta = r \in \R$), the per-oscillator equilibrium satisfies $(Kr/2)z^2 + i\omega z - Kr/2 = 0$. For locked oscillators ($|\omega| < Kr$):

\begin{proposition}[\texttt{lockedEquil\_is\_equil}, \texttt{lockedEquil\_normSq}, \texttt{lockedEquil\_stability\_neg}; 0 sorry]\label{prop:locked-equil}
The locked equilibrium
\[
z^*(\omega) = \frac{\sqrt{K^2r^2 - \omega^2} - i\omega}{Kr}
\]
satisfies: (i)~$z^*$ solves $\dot z = 0$, (ii)~$|z^*| = 1$, (iii)~$0 < \re(z^*) \leq 1$ with equality iff $\omega = 0$, (iv)~the linearized eigenvalue $\lambda = -i\omega - Kr\,z^*$ has $\re(\lambda) = -\sqrt{K^2r^2 - \omega^2} < 0$, (v)~$\re(z^*)$ is monotone increasing in~$r$, and (vi)~the self-consistency function $F(r) := \int_{|\omega|<Kr} \re(z^*(\omega))\,g(\omega)\,\mathrm{d}\omega$ satisfies $F(0) = 0$, $F \geq 0$, $F \leq \int g$, and $F$ is monotone.
\end{proposition}
\begin{proof}
Direct substitution: $(Kr/2)(z^{*2} - 1) + i\omega z^* = 0$ reduces to the identity $D^2 + \omega^2 = K^2r^2$ where $D = \sqrt{K^2r^2-\omega^2}$. Parts (ii)--(iv) follow from the explicit formula. Part (v) uses $\re(z^*)^2 = 1 - \omega^2/(K^2r^2)$, which is increasing in~$r$. Part (vi): $F(0) = 0$ since the locked domain is empty, monotonicity follows from (v) plus the expanding locked domain (\texttt{ComplexOALockedEquil.lean}, 270 lines, 19 theorems, 0~sorry).
\end{proof}

\subsection{V\texorpdfstring{$'$}{\'} decomposition}

By Proposition~\ref{prop:rotation}, the rotation $-i\omega(z-z^*)$ contributes zero to~$V'$. The remaining coupling terms decompose as follows.

\begin{proposition}[\texttt{complex\_V\_deriv\_eq\_QDS}, 0 sorry]\label{prop:v-prime}
Define $Q_c := \int |z-z^*|^2\,\re(z+z^*)\,g$, $S_c := \int \re(\bar d(1-z^2))\,g$, and $D_c := \int \re(\bar d)\,g$ where $d = z - z^*$. If $r(t) - r^* = D_c$ (order parameter identity), then
\[
V'(t) = K\Bigl(-\frac{r^*}{2}\,Q_c + \frac{D_c}{2}\,S_c\Bigr).
\]
\end{proposition}
\begin{proof}
Factor the per-$\omega$ integrand (\texttt{complexVDerivIntegrand\_expand}, by \texttt{ring}), then apply linearity of integration with the order parameter substitution $r(t) - r^* = D_c$. Integrability of the $S_c$ and $Q_c$ integrands is required as hypothesis (\texttt{ComplexPairBoundProof.lean}, 0~sorry).
\end{proof}

By the real-case pair bound architecture, $V' \leq 0$ would follow from $r^*Q_c - D_cS_c \geq 0$, which is the Fubini double-integral nonnegativity. The pointwise pair bound is provably \emph{false} for complex~$z$ (\texttt{complex\_pair\_pointwise\_false}), so the integral nonnegativity requires genuine Fubini analysis; this remains open.

\subsection{Formal order-parameter bridge}\label{sec:bridge}

\begin{lemma}[Formal order-parameter bridge, 0 sorry]\label{lem:bridge}
Let $R_{\mathrm{OA}}, R_{\mathrm{full}} : [0,\infty) \to \C$ and let $r_* > 0$. If
\[
|R_{\mathrm{OA}}(t)| \to r_*
\qquad\text{and}\qquad
|R_{\mathrm{full}}(t) - R_{\mathrm{OA}}(t)| \to 0,
\]
then $|R_{\mathrm{full}}(t)| \to r_*$.
\end{lemma}
\begin{proof}
By the reverse triangle inequality,
\[
\bigl||R_{\mathrm{full}}(t)| - |R_{\mathrm{OA}}(t)|\bigr| \leq |R_{\mathrm{full}}(t) - R_{\mathrm{OA}}(t)|.
\]
The second hypothesis makes the right-hand side tend to zero, while the first gives $|R_{\mathrm{OA}}(t)| \to r_*$. Hence $|R_{\mathrm{full}}(t)| \to r_*$ (\texttt{squeeze\_zero\_norm + abs\_norm\_sub\_norm\_le}, 0~sorry).
\end{proof}

\begin{remark}[Status of the bridge lemma]\label{rem:bridge-hypotheses}
Lemma~\ref{lem:bridge} is a formal consequence of the reverse triangle inequality. It is not, by itself, a proof of full Kuramoto PDE stability. Its two assumptions---convergence of the OA order parameter and asymptotic closeness of the full PDE order parameter to the OA order parameter---are \emph{external hypotheses that are not proved in this manuscript}. We describe below what would be needed to discharge them.

\medskip\noindent\textbf{Reference~\cite{dietert2017}: nonlinear Landau damping.} Dietert~\cite{dietert2017} (building on~\cite{dietert-thesis}) proves local nonlinear Landau damping / orbital stability for linearly stable partially locked states of the full Kuramoto PDE, after removing the neutral rotation mode. In~\cite{dietert2017}, the decaying quantity $\eta(t)$ is the perturbation order parameter $u_1(t,0)$ in a rotating frame, not automatically the full physical order parameter and not an OA order parameter. The hypotheses include Sobolev regularity of the perturbation, smallness near a linearly stable PLS, and phase/gauge fixing.

\medskip\noindent\textbf{Reference~\cite{dietert-fernandez2018}: OA manifold attractivity.} Dietert--Fernandez~\cite{dietert-fernandez2018} controls attraction to the OA manifold through Fourier-side OA defects
\[
w_{n,m} := \hat f_{n+m} * \hat g - \hat f_n * \hat f_m,
\]
which vanish on the OA manifold. Their size is measured in a weighted Fourier--Sobolev norm $\|w\|_{H^1_{e^{a\tau}}}$; exponential decay $\|w(t)\|_{H^1_{e^{a\tau}}} \leq \|w(0)\|_{H^1_{e^{a\tau}}}\,e^{-at}$ is proved in this topology. This is not the same as directly proving the existence of an OA trajectory satisfying $|R_{\mathrm{full}}(t) - R_{\mathrm{OA}}(t)| \to 0$: the defect norm is a global Fourier--Sobolev quantity, and extracting a scalar order-parameter bound from it requires additional trace/evaluation estimates.

\medskip\noindent\textbf{Matching gap.} To derive the two scalar assumptions of Lemma~\ref{lem:bridge} from~\cite{dietert2017} and~\cite{dietert-fernandez2018}, one would still need: choosing the correct OA representative, matching phase conventions and rotating frames, proving trace/evaluation estimates from the Fourier--Sobolev defect norm to the scalar order parameter, and verifying all local smallness, regularity, and linear stability hypotheses. This matching is not carried out here. The corresponding Lean bridge declaration (0~sorry) takes OA convergence as a regular hypothesis and encodes PDE-to-OA closeness via one \texttt{opaque} type with one \texttt{axiom}; the two convergence inputs remain external assumptions whose hypotheses are not matched in this paper.
\end{remark}

% ============================================================
\section{Discussion}\label{sec:discussion}
% ============================================================

\subsection{Four tiers of formalization}

The results are organized into four tiers of decreasing certainty:

\medskip
\noindent\textbf{Theorem A: Real scalar complete} (\texttt{real\_scalar\_complete}). Fully formalized, 0~sorry, 0~axioms. Basin stability for the scalar dissipative model~\eqref{eq:oa} with finite first moment $\int \gamma\,g < \infty$. When $\gamma(\omega) \lesssim 1 + |\omega|$: Gaussian, Student-$t$ ($\nu > 1$), compact support, mixtures.

\medskip
\noindent\textbf{Theorem B: Lorentzian} (\texttt{lorentzian\_ode\_global\_stability\_complete}). Fully formalized, 0~sorry, 0~axioms. Global stability for the Bernoulli ODE with explicit solution formula. Finite-dimensional, no measure theory needed.

\medskip
\noindent\textbf{Theorem C: $n$-pole} (\texttt{kuramoto\_solved}). Fully formalized, 0~sorry, 0~axioms. Exponential decay rate $\mu = K\delta\delta^*$. Requires bounded $\gamma \leq \gamma_{\max}$, uniform persistence $\delta \leq \alpha(\omega,t)$ for all $\omega,t$, and a pre-existing equilibrium $\alpha^*$.

\medskip
\noindent\textbf{Lemma D: Conditional bridge.} Formal bridge lemma (reverse triangle inequality), 0~sorry. OA convergence is a regular hypothesis; PDE-to-OA closeness is encoded via one \texttt{opaque} type and one \texttt{axiom}. The two convergence inputs are external assumptions not derived from~\cite{dietert2017} or~\cite{dietert-fernandez2018} here; see Remark~\ref{rem:bridge-hypotheses}.

\subsection{Theorem table}

\begin{center}
\small
\begin{tabular}{llllll}
\textbf{Name} & \textbf{Model} & \textbf{Hypotheses} & \textbf{Conclusion} & \textbf{Sorry} & \textbf{Axioms} \\
\hline
\texttt{real\_scalar\_complete} & Scalar dissipative & $\int\gamma\,g < \infty$, $V(0) < (r^*)^2$ & $r(t) \to r^*$ & 0 & 0 \\
\texttt{lorentzian\_ode\_global\_stability\_complete} & Lorentzian ODE & $K > 2\gamma$, $r_0 \in (0,1)$ & $r(t) \to r^*$ & 0 & 0 \\
\texttt{kuramoto\_solved} & Bounded $\gamma$ & $\gamma \leq \gamma_{\max}$, uniform persistence & $r(t) \to r^*$ & 0 & 0 \\
(bridge declaration) & Formal conditional bridge & OA conv.\ + PDE-to-OA closeness & $|R_{\mathrm{full}}| \to r_*$ & 0 & 1 \\
\end{tabular}
\end{center}

\subsection{Novelty}

This work contributes four things that are new:

\begin{enumerate}
\item \textbf{First machine-checked Kuramoto result.} To our knowledge, this is the first formalization in any proof assistant (Lean, Coq, Isabelle) addressing any aspect of the Kuramoto model.

\item \textbf{Novel proof technique.} The pair bound (Proposition~\ref{prop:pair})---a pointwise algebraic inequality lifted to the continuum double integral via Fubini---has no precedent in the Kuramoto literature. Prior approaches used Fourier analysis \cite{dietert2016}, Volterra equations \cite{dietert-fernandez2018}, or spectral theory \cite{chiba2015}.

\item \textbf{Rotation cancellation.} The rotation term $-i\omega$ in the complex OA cancels in the $V$~derivative after equilibrium subtraction (Proposition~\ref{prop:rotation}). This removes the rotation contribution from $V'$, but the remaining nonlinear coupling terms differ structurally from the scalar case. Convergence of the complex OA is not proved here (see Remark~\ref{rem:bridge-hypotheses}).

\item \textbf{Finite first moment coverage.} Prior rigorous OA analysis was limited to Lorentzian/rational distributions. The finite first moment condition $\int \gamma\,g\,\mathrm{d}\omega < \infty$ covers Student-$t$ with $\nu > 1$ (including the infinite-variance case $1 < \nu \leq 2$), Gaussian, compact support, and all mixtures.
\end{enumerate}

\subsection{Limitations}

\begin{enumerate}
\item \textbf{Basin of attraction.} The continuum result (Theorem~A) requires $V(0) < (r^*)^2$, an explicit but \emph{local} condition. True global stability (any $r(0) > 0$) is proved only for the Lorentzian ODE (Theorem~B). Theorem~C requires uniform persistence as a hypothesis.

\item \textbf{Unmatched external assumptions.} Lemma~D is only the reverse triangle inequality. Its two mathematical assumptions---OA convergence and PDE-to-OA closeness---are not derived from~\cite{dietert2017} and~\cite{dietert-fernandez2018} here. The matching (phase conventions, rotating frames, trace estimates, smallness/regularity verification) is not carried out. In Lean, OA convergence is a regular hypothesis; PDE-to-OA closeness is encoded via one \texttt{opaque} type and one \texttt{axiom}. See Remark~\ref{rem:bridge-hypotheses} for details.

\item \textbf{Conclusion is $\re(\eta)^2$, not $|\eta|^2$.} On the symmetric subspace $\eta \in \R$ (proved, 0~sorry), so $\re(\eta)^2 = |\eta|^2$. But the Lean type states $\re(\eta)^2 \to (r^*)^2$ rather than $|\eta|^2 \to (r^*)^2$.

\item \textbf{Scalar model scope.} The scalar dissipative model~\eqref{eq:oa} is a separate dynamical system from the standard complex OA~\eqref{eq:complex-oa}.
\end{enumerate}

\subsection{Open problems}

\begin{enumerate}
\item Remove the $V(0) < (r^*)^2$ basin condition (global stability for any $r(0) > 0$). The locked equilibrium structure (Proposition~\ref{prop:locked-equil}) and self-consistency interface are formalized; the remaining gap is proving $\liminf r(t) > 0$ for $K > K_c$ via instability of incoherence.
\item Match and formalize the nonlinear Landau damping result of~\cite{dietert2017}: identify the correct rotating-frame order parameter, verify linear stability and Sobolev smallness hypotheses, and extract a scalar convergence bound in the physical frame.
\item Match and formalize the OA attractivity result of~\cite{dietert-fernandez2018}: extract from the Fourier--Sobolev defect norm a scalar bound $|R_{\mathrm{full}}(t) - R_{\mathrm{OA}}(t)| \to 0$ via trace/evaluation estimates.
\item Extend to distributions without finite first moment (Student-$t$ with $\nu \leq 1$) beyond the Lorentzian special case.
\end{enumerate}

\begin{thebibliography}{10}
\bibitem{dietert2016} H.~Dietert, Stability and bifurcation for the Kuramoto model, \emph{J. Math. Pures Appl.}~\textbf{105} (2016), 451--489.
\bibitem{dietert-fernandez2018} H.~Dietert and B.~Fernandez, The mathematics of asymptotic stability in the Kuramoto model, \emph{Proc. R. Soc. A}~\textbf{474} (2018), 20180467.
\bibitem{dietert-thesis} H.~Dietert, \emph{Stability of the Kuramoto model}, PhD thesis, Universit\'e Paris Diderot, 2016.
\bibitem{kuramoto1975} Y.~Kuramoto, Self-entrainment of a population of coupled nonlinear oscillators, in \emph{International Symposium on Mathematical Problems in Theoretical Physics}, Lecture Notes in Physics~\textbf{39}, Springer, 1975, pp.~420--422.
\bibitem{ott-antonsen2008} E.~Ott and T.~M.~Antonsen, Low dimensional behavior of large systems of globally coupled oscillators, \emph{Chaos}~\textbf{18} (2008), 037113.
\bibitem{chiba2015} H.~Chiba, A proof of the Kuramoto conjecture for a bifurcation structure of the infinite-dimensional Kuramoto model, \emph{Ergod. Th. \& Dynam. Sys.}~\textbf{35} (2015), 762--834.
\bibitem{engelbrecht2020} J.~R.~Engelbrecht and R.~Mirollo, Is the Ott-Antonsen manifold attracting?, \emph{Phys. Rev. Research}~\textbf{2} (2020), 023057.
\bibitem{strogatz2000} S.~H.~Strogatz, From Kuramoto to Crawford: exploring the onset of synchronization in populations of coupled oscillators, \emph{Physica D}~\textbf{143} (2000), 1--20.
\bibitem{dietert2017} H.~Dietert, Stability of partially locked states in the Kuramoto model through Landau damping with Sobolev regularity, \emph{arXiv:1707.03475} (2017, revised 2020).
\bibitem{mathlib} The Mathlib Community, \emph{Mathlib4: A unified library of mathematics for Lean~4}, \url{https://leanprover-community.github.io/mathlib4_docs/} (2024).
\end{thebibliography}

\end{document}
